\chapter{Background and Related Work}

This chapter follows one line of reasoning rather than surveying bodies of literature side by side.
Is the problem real, what has already been tried, has removing things instead of adding them been
tried and did it work, can any of this be built on hardware people actually own, and how should the
result be tested. The review is selective on that basis: it includes work that establishes the
problem, backs a value this system uses, or supplies a method the evaluation adopts, and it leaves
work that merely neighbours the topic to the surveys that cover it.

\section{Why the Periphery Costs Something}

The computational account of visual attention underpinning this dissertation is the salience-map
model of \citet{itti1998}: attention is drawn to locations whose centre--surround contrast, in
intensity, colour opponency, or orientation, differs most from their neighbourhood. Three working
assumptions, adopted from this model, recur throughout the design. First, the model treats
\textit{contrast as the master variable}, so any
manipulation raising local contrast should attract gaze and any manipulation flattening it should
repel gaze.
Second, salience is channel-specific and the channels are not equal: in every head-to-head comparison
in the guidance literature, luminance manipulations outperform chromatic ones
\citep{bailey2009,grogorick2017}. This is why every mode in this system dims before it does anything
else. Third, salience operates \textit{pre-attentively}, so capture by a high-contrast peripheral
event happens before, and regardless of, the observer's intentions. On this account a periphery
whose salience has been flattened should stop issuing capture events, which is the mechanism the
dark modes aim to exploit.
\citet{duan2022} extend the account to augmented reality and show that the opacity of AR content
directly modulates where attention lands, which makes opacity a dose rather than a style and the
attenuation axis of Chapter 3 an axis with a middle rather than an on/off switch.

Whereas salience explains \textit{how} the periphery captures attention, Cognitive Load Theory explains
why that capture is costly. Extraneous load, the processing demanded by information irrelevant to the
task, competes for the same limited working-memory resources as the task itself. \citet{buchner2022}
review the evidence that AR can reduce, not only add, cognitive load when used to structure or filter
what the user must process. This is the theoretical justification for the whole project: peripheral
diminishment is an extraneous-load intervention. It does not make the user smarter or the task
easier. It removes competing processing demands at the perceptual source.

That framing also identifies the project's central risk, which is why this dissertation pairs a
benefit hypothesis with a safety hypothesis rather than testing benefit alone. Attention capture by
peripheral events is not a design flaw of the human visual system. It is a safety mechanism.
Suppressing it trades focus against situational awareness, a tension documented empirically in the
DR literature \citep{murph2021,mclaughlin2025} and formalised in Chapter 5 as an equivalence-tested
safety hypothesis rather than an afterthought. The evaluation is designed to measure benefit and
cost together.

What was missing until recently was the distractor-cost baseline in immersive settings.
\citet[N = 66]{ai2025} supply it: peripheral visual distractors in a virtual classroom significantly
increase both commission errors (1.33 $\rightarrow$ 3.15) and omission errors on a Go/No-go
continuous performance task, while leaving mean reaction time unchanged. That single result does
double duty here: it establishes that \textit{errors, not latency}, are the sensitive currency of
distraction cost, fixing the primary dependent variable of Chapter 5, and it quantifies the effect
the vignette must reduce for the primary hypothesis to be supported.

\section{Guiding Attention by Adding}

The techniques for steering attention in a mediated view form a spectrum from imperceptible to
overt, and what matters here is what the additive approach costs.

The subtle gaze direction family originates with \citet{bailey2009}: a brief luminance modulation in
the viewer's periphery attracts a saccade, and the cue is terminated the moment gaze moves toward it,
so the viewer never foveates the modulation and remains unaware of being steered. Luminance
modulation outperformed warm--cool chromatic modulation, an early instance of the luminance-dominance
regularity described in Section 2.1. Two results from this family set hard constraints on the present design.
\citet{grogorick2018} compared five guidance methods in an immersive dome with N = 102: local
magnification achieved the best guidance-per-noticeability ratio, while whole-periphery blur was the
least subtle and actively irritated 59 of 102 participants. Their conclusion that no method remains
completely unnoticed, together with the blur-irritation finding, decided two things here.
Whole-periphery blur is not carried forward as an evaluated mode, and the design accepts
noticeability rather than chasing imperceptibility. Second, \citet{erickson2023} found hue shifts the
most confusing of hue, saturation and value, which is why every mode manipulates saturation and
luminance and none shifts hue.

Previous work suggests slow ramps can escape the noticeability trade-off. \citet{hata2016} showed
that a visual change introduced gradually enough stays below the awareness threshold while still
biasing gaze, because the guidance threshold sits below the awareness threshold. This is the perceptual
justification for the system's formation animation, where effects ramp in over seconds rather than
switching on.

Where subtle gaze direction adds a transient stimulus, saliency \textit{modulation} re-grades the
image itself. This is the direct technical lineage of this system's colour- and saliency-gated
re-grading mode (ColorPop). The foundational demonstration is \citet{veas2011}, who raised salience in a focus region and
suppressed it in context regions of video imagery using an Itti-style conspicuity model. Modulated
video was statistically indistinguishable from unmodulated video for naive viewers yet produced
faster first fixations on target regions, which remains the strongest evidence that video-mediated
reality can be re-graded imperceptibly with measurable attentional consequences. A passthrough
headset is exactly video mediation. \citet{kalkofen2007} contribute the focus-and-context vocabulary
used throughout: keep a focus region maximally legible while context is compressed but not
destroyed.

The modern anchor is \citet{sutton2022}, saliency modulation of the \textit{real world} viewed
through AR glasses, implemented as saturation adjustment plus sigmoidal contrast remapping and
deliberately avoiding hue shifts. Modulation drew significantly more fixations to target regions but,
an honest and important negative, was \textit{not} imperceptible. Against an overt circle cue,
modulation's advantage was that it preserved natural scene exploration where explicit markers glue
gaze to themselves. This dissertation's ColorPop implements Sutton et al.'s exact sigmoidal contrast
recipe on its salient-colour channel, inheriting the philosophy: accept visibility, preserve
exploration, re-grade rather than annotate. The critical difference is platform. Sutton et al. built
a bench-mounted optical see-through rig whose additive optics can only add light. The present system
runs on consumer video see-through hardware with full subtractive control of its rendered layers.

An eye-tracked comparison in cinematic VR speaks to which \textit{diminishment} channel works.
Comparing area darkening, context-based darkening and desaturation for steering attention in
360$^{\circ}$ video, area darkening guided attention most effectively while desaturation alone had
almost no guidance effect \citep{wang2020}. It is a single result from one small-venue study rather
than a consensus finding, but it corroborates this project's mode hierarchy: the dark modes are the
guidance workhorses, and desaturation appears only as one ingredient inside ColorPop's compound
re-grading, never as a mode of its own.

At the overt end, the Attention Funnel \citep{biocca2006} establishes that aggressive guidance
pays in speed what it costs in scene engagement. Its failure mode is \textit{attentional
tunnelling}, excessive allocation of attention to guided content at the cost of neglecting the
surroundings, isolated experimentally as a function of task demand \citep{syiem2021}. That is the
risk most relevant to Hard Dark, and the reason its evaluation measures peripheral awareness rather
than assuming it.

\section{Guiding Attention by Removing}

\citet{mori2017} provide the canonical survey and taxonomy of diminished reality, the family of
techniques that conceal, eliminate, or see through real objects in a mediated view. Their survey
establishes both the field's legitimacy and its historical centre of gravity: most classical DR work
concerns object removal, inpainting a region so the object appears gone, which is computationally
demanding and semantically targeted. The present work sits in a different corner of that taxonomy. It
does not remove objects but \textit{attenuates regions}, trading semantic precision for robustness,
generality and real-time feasibility on mobile hardware.

\citet{cheng2022} supply the field's most direct study of how users \textit{want} reality diminished.
Across two studies (N = 16 and N = 12, conducted inside VR simulations of DR because no deployable DR
hardware existed) they found participants preferred partial opacity reduction over complete removal,
and wanted context-preserving diminishment, outlines or residual structure indicating that something
is there even when its detail is suppressed. Both findings are load-bearing. The preference for partial attenuation fixes the opacity ceiling of
this system's partial peripheral occlusion mode (Soft Dark), and the context-preservation finding
motivates treating Hard Dark's total blackout as the \textit{extreme} of the design space rather
than its default. Equally load-bearing is their method: they could only simulate DR inside VR, where this
dissertation runs on the real passthrough hardware their participants could only imagine.

The most recent conceptual frame for this family is \textit{visual noise cancellation}: by analogy
with acoustic noise cancellation, a head-worn display actively moderates the visual environment,
attenuating noise while passing signal \citep{hong2024}. The analogy is productive because it imports
an architecture. Acoustic cancellation requires a microphone, a processing stage, and a speaker,
which is precisely the camera $\rightarrow$ shader $\rightarrow$ display pipeline of a
video-see-through headset. Against the prior art surfaced by the structured search of Section 2.5.1,
no software-defined visual noise cancellation system on consumer video-passthrough hardware was
found, and this dissertation can be read as one, where each mode is a cancellation profile and the
profiles differ in what they classify as noise. The analogy also imports a useful discipline
from its original: noise-cancelling headphones are evaluated on measured attenuation \textit{and} on
what they cost the wearer in awareness, which is exactly the hypothesis pairing of Chapter 5.

The empirical case that diminishment helps focus is recent but consistent. \citet{mclaughlin2025}
provide the closest published analogue to this dissertation's primary hypothesis: across two
within-subject experiments, universal DR-style attenuation of distractors, the condition structurally
closest to this project's peripheral vignette, produced fewer errors ($\beta = -0.37$, $p = .021$) and
lower subjective workload than context-aware removal. Their trial-level mixed-effects analysis is
adopted directly by this dissertation's analysis plan, and their measurement of awareness of off-task
objects, finding the focus/awareness trade-off real but manageable, is the empirical ancestor of this
work's safety hypothesis. \citet{lee2025} approach the same goal from the object side: AR-camouflaging a single distracting
object restored working-memory performance to a level statistically comparable with physically
removing it (N = 60). \citet{murph2021} document the same benefit--risk structure in a medical
assembly context, where DR lowered cognitive workload but risked situational-awareness loss, and in
later work describe gradually reintroducing distractions as tolerance grows \citep{murph2022}, which
is the source of this system's temporal fade transitions.

Two 2026 threads confirm both the demand for DR-for-attention and its immaturity. A co-design
study with fifteen students with ADHD developed the concept of attentional diminishment but built no
working prototype \citep{exploring}, and a CHI 2026 study on Apple Vision Pro obscured specific
individuals to alleviate social discomfort, which is overlay-based diminishment on consumer hardware
but targeted occlusion rather than peripheral guidance \citep{obscuring}. Neither implements
peripheral diminishment for attention on real passthrough. Further back, FocalSpace \citep{yao2013}
demonstrated the 2D ancestor of this project's workstation scenario: synthetic background blur in
video conferencing, improving memory for meeting content.

\section{What the Periphery Tolerates}

The perception literature describes peripheral vision not as blurry central vision but as a
differently tuned system, one that notices contrast loss and motion change most, and chromatic
detail loss least. One perceptual finding is
directly responsible for a parameter in this system: \citet{patney2016} established that peripheral
blur is detected primarily through the \textit{contrast reduction} it causes, and that restoring
local contrast after blurring roughly doubles the tolerable blur radius. That chain also offers a
mechanistic explanation for the irritation finding of \citet{grogorick2018}, since naive
whole-periphery blur announces itself through contrast collapse. It is why the dark modes attenuate
luminance, which the periphery tolerates better, rather than blurring structure.

A second literature approaches the periphery from the comfort side, and the manner of restriction
matters more than its degree. \citet{fernandes2016} showed a soft, dynamically applied field-of-view
vignette reduces discomfort in VR locomotion without measurable presence loss. But
\citet{norouzi2018} found that vignetting \textit{coupled to the user's own head motion} actually
\textbf{increased} sickness relative to no vignette, a directly actionable negative result that
justifies this system's inverted motion policy, where effects fade \textit{out} during head motion
and restore during stability. \citet{cao2021} found granulated rest frames outperform opaque black
restrictors on search performance while preserving peripheral awareness, the source for this
system's world-locked static-grain variant (GranulatedPeriphery). \citet{barhorstcates2016} contribute the sobering datum at the far end:
spatial learning survives restriction down to a 10$^{\circ}$ field, but anxiety is elevated in
\textit{all} restricted conditions and rises monotonically as the field narrows. \citet{waldner2014}
map the comfort envelope for the flicker channel, giving the parameters this project adopts for its
detection-highlighting pulse rather than anything in the photosensitivity band.

Collectively this draws the corridor the system must fly through. Peripheral attenuation is tolerable
and even beneficial \textit{if} it is soft-edged, decoupled from head motion, and mindful that total
blackout carries a measurable anxiety cost. Those constraints are visible in the final system as the
dark modes' graduated opacity ceilings, every mode's motion-triggered suppression, and the
evaluation's insistence on measuring peripheral awareness alongside benefit. One honest gap remains,
inherited rather than resolved: no study measures the comfort of peripheral \textit{desaturation}
specifically over sessions longer than ten minutes.

\section{What Consumer Hardware Permits}

Everything above presumes the ability to manipulate what the user sees of reality. On consumer
hardware in 2026 that ability is narrowly rationed, and Chapter 3 develops the rationing in full as a
three-tier design space. Two literature points attach to it here. The one published feasibility study
of on-device camera compositing projects 720p30 with thermal throttling within five to ten minutes, a
simulation-based estimate rather than a measured run, which functions as a warning about sustained
camera processing \citep{laghari2025}. And Varjo's research headsets grant full camera control, yet
this review found no published study that shader-processes their passthrough region-selectively for
attention. The question appears to be unexplored so far, although work not surfaced by the search
may exist.

One further platform datum constrains not the system but its \textit{evaluation}: psychophysical
measurement shows that no current video see-through headset reaches normal human acuity at any light
level, with Quest 3 significantly worse in low light \citep{wang2026}. Any evaluation task requiring
fine real-world detail read through passthrough therefore confounds the manipulation under test with
the medium's acuity ceiling. This drives two decisions in Chapter 5: the focal task uses large
look-alike shapes whose legibility through passthrough is verified by a pilot gate, rather than
anything requiring fine detail, and room lighting is bright and constant across sessions.

\subsection{Prior systems, and a documented search}
\label{sec:priorart}

Every constituent mechanic of this architecture has a documented precedent, and none of the
precedents combine them. QuestCameraKit demonstrates camera-fed materials with GPU shaders
\citep{coviello2025}, but always as content \textit{within} the scene rather than a full-surround
treated periphery. Meta's Passthrough Transitioning motif uses a head-centred sphere whose
transparency reveals system passthrough \citep{meta2024c}, but the sphere carries virtual content for
scene fades rather than a processed camera feed. Alpha punch-through to the passthrough underlay is
standard technique \citep{meta2024b}, shipped commercially by Immersed's Passthrough Portals
\citep{immersed2022}, but the hole is always surrounded by \textit{virtual} workspace rather than by
re-rendered reality. A peripheral dimming vignette shipped in Quest OS v67, but explicitly not over
passthrough \citep{meta2024d}. Two patents describe adjacent ideas with no evidence either was
implemented \citep{patents}.

A structured prior-art search conducted for this dissertation on 5 July 2026, comprising 26 web
queries and 4 GitHub API queries across ISMAR, CHI, UIST, IEEE VR, DIS and arXiv from 2022 to 2026,
Meta developer documentation, all 136 public forks of the PCA samples repository, XR trade press, and
the visionOS and Varjo API surfaces, found \textbf{no publication, product, or open-source project
that composites system passthrough with a live shader-processed camera feed in a single view, for any
purpose}, and none that implements peripheral passthrough manipulation for attention guidance on a
standalone consumer headset. The full query protocol is archived in the project record. The nearest
systems are the five named above, each contributing one ingredient. Chapter 3, Section 3.5 states the
resulting novelty claim, scoped precisely to the composite rather than to DR on passthrough in
general.

\section{Evaluating Attention Guidance in XR}

The evaluation norms of this field are within-subject designs with modest samples and many
repeated measures. \citet{caine2016}'s census of CHI found the modal sample size to be twelve, with
in-person studies drawing their power from paired comparisons and trial-level repetition rather than
headcount. For non-normal aggregates the HCI norm has been the aligned rank transform
\citep{wobbrock2011}, though recent critique of its error properties \citep{tsandilas2024} motivates
this dissertation's preference for mixed models with Wilcoxon fall-backs.

Quest 3 carries no eye tracker, so the measures must come from behaviour and head telemetry, and the
literature validates both channels. Reaction-time probe detection is standardised as the Detection
Response Task \citep{iso2016}, and the Peripheral Detection Task places its probes at
11--23$^{\circ}$ eccentricity \citep{jahn2005}, just outside the roughly 8$^{\circ}$ road-centre
region where driver gaze concentrates \citep{victor2005}. That lineage fixes the eccentricity bands
of Chapter 5. Head orientation is a validated proxy for gaze at the eccentricities that matter here:
\citet{higgins2022} validated head pose as a gaze surrogate for object-level attribution in VR, and
\citet{sitzmann2018} showed with concurrent head and eye tracking that fixations concentrate at low
head velocity. Simulator sickness is measured with the nine-item VRSQ \citep{kim2018}, administered
before \textit{and} after exposure because post-only scores are uninterpretable without a baseline
\citep{brown2022}. The study administers no per-condition workload instrument, and Chapter 6 records
what that costs. Finally, equivalence
claims of the ``no meaningful loss of peripheral awareness'' kind require dedicated machinery: the
two-one-sided-tests procedure with a pre-registered margin \citep{lakens2017}, which converts a
would-be null result into a falsifiable claim.

\section{Gaps and Position}

Read together, the literature above leaves five identifiable gaps, three of which this dissertation
addresses in whole or part. Table~\ref{tab:gaps} states each gap and what this work does about it.

\begin{longtable}{p{0.05\textwidth}|p{0.35\textwidth}|p{0.50\textwidth}}
\caption{Research gaps identified by the review, and the standing of each in this dissertation.}
\label{tab:gaps} \\
\hline
\textbf{\#} & \textbf{Gap} & \textbf{Status in this dissertation} \\
\hline
\endfirsthead
\hline
\textbf{\#} & \textbf{Gap} & \textbf{Status in this dissertation} \\
\hline
\endhead
1 & \textbf{No unified multi-mode DR system.} Guidance techniques and diminishment effects have been studied singly. No prior work combines multiple DR guidance modes in one configurable system for comparison within one design. & Addressed: ten modes built in one configurable system spanning the design-space tiers of Chapter 3, of which two are carried to confirmatory comparison within one study. \\
\hline
2 & \textbf{No hybrid of system passthrough and app-processed camera feed in one view.} Colour-only global restyling of system passthrough exists. Full pixel control of an inferior camera feed exists. No prior work composites the two, and none implements peripheral passthrough manipulation for attention on a standalone consumer HMD (documented search, Section~\ref{sec:priorart}). & Addressed: the window-as-absence hybrid architecture of Chapters 3 and 4, the dissertation's strongest and most precisely scoped novelty claim. \\
\hline
3 & \textbf{No temporal transition mechanics.} Gradual formation, easing, and motion-contingent suppression of DR effects have not been formally parameterised or studied, despite adjacent evidence that ramps evade noticeability \citep{hata2016} and that head-coupled strengthening harms comfort \citep{norouzi2018}. & Partially addressed: formation ramps and inverted motion-coupling are implemented and parameterised in Chapter 4 and grounded in the cited evidence. Their isolated evaluation remains future work. \\
\hline
4 & \textbf{No user-controlled DR intensity.} User-initiated, graded release and re-application of diminishment is unstudied. & Not addressed, groundwork only. The trigger-painted window gives users spatial but not intensity control. \\
\hline
5 & \textbf{No dual-channel DR technique.} Desaturation and blur have been studied separately, never combined and parameterised as a compound diminishment channel in AR attention guidance. & Not addressed as a confirmatory question. The compound channel exists inside ColorPop, but the study tests the unconfounded dark modes confirmatorily and leaves channel decomposition to future work. \\
\hline
\end{longtable}

Two published studies stand closest to this dissertation, and the delta from each defines its
contribution. \textbf{\citet{cheng2022}} established what users want from diminishment, but inside VR
simulations, because deployable DR did not exist. Theirs is a perception-and-preference study of
imagined systems. \textbf{\citet{mclaughlin2025}} established that distractor attenuation improves
task performance, but likewise inside a virtual task environment, with the diminishment simulated by
the VR scene itself. Both are simulation studies \textit{of necessity}. This dissertation contributes
the piece both were missing: a working diminishment system on real consumer passthrough hardware,
evaluated with a design that pairs a real-hardware block against a simulation block within the same
participants. If the two blocks agree in direction, the simulation methodology of Cheng and
McLaughlin gains a real-hardware anchor. If they diverge, the divergence localises what simulation
misses. Either way the evaluation assembles the field's strongest methodological norms, for the
first time, around a deployable diminished-reality attention system.

One further applied result matters to that evaluation, because it uses driving \textit{footage} as
a stimulus context. AR cueing of roadside hazards improves response time \textit{without} degrading
detection of non-target objects \citep[N = 27]{rusch2013}, an early demonstration that guidance
benefit and situational-awareness cost are separable, which is the separation this dissertation's
secondary hypotheses formalise. That literature motivates the scenario aesthetics and supplies the
salience-engineering precedent for ColorPop's red, orange and green hierarchy. It licenses nothing
further: per Chapter 1, no claim about driving, drivers, or road safety is made anywhere in this
dissertation.
